\chapter{Detector calibration and the zero-lag access boundary}
\label{ch:ci-calibration}

The size of an observed fluctuation depends on the detector as well as on the
physical process. Calibrating that dependence can reveal irreversibility in a
positive, monotonically decreasing correlation, before a negative correlation
or revival appears. The same analysis also identifies a limit: arbitrarily
precise measurements at a fixed positive-lag pair menu need not determine the
calibration quantity itself.

\section{The conditional second moment}
Let a fixed, noninvasive binary detector have conditional mean
$f_i=\mathbb E[Y\mid X=i]\in[-1,1]$. Under the stationary law $\pi$, define
\begin{equation}
 m=\sum_i\pi_i f_i,\qquad A=\sum_i\pi_i f_i^2,
 \qquad u=m^2,\qquad0\leq u\leq A\leq1.
 \label{eq:ci-mA}
\end{equation}
Two readings with conditionally independent noise at the same unchanged hidden
state obey
\begin{equation}
 \mathbb E[Y^{(1)}Y^{(2)}]=A,
 \qquad\Pr(Y^{(1)}=Y^{(2)})=\frac{1+A}{2}.
 \label{eq:ci-zero-lag}
\end{equation}
Squaring one binary reading instead gives $Y^2=1$ and does not measure $A$.
The two readings are generally dependent through their common hidden state.

For an instantaneous detector, $m$ and $A$ agree under $Q$ and its additive
reversibilization because $\pi$ and $f$ are unchanged. This transfer need not
hold for nonlinear finite-exposure path functionals. A deterministic binary
state detector already has $A=1$. If $f(X)$ itself is reliably observed, its
square directly measures $A$; this is additional access, outside the binary-only
obstruction proved below.

\section{Bounded analog observations supply an upper calibration}
A trusted instantaneous analog record can improve the comparison class even
without a same-state replicate. Let $Z\in[-1,1]$ be the output of a fixed,
noninvasive memoryless channel, and write $f_i=\mathbb E[Z\mid X=i]$. Then
\begin{equation}
 m=\mathbb E Z,\qquad
 A=\mathbb E f(X)^2\leq B:=\mathbb E Z^2,
 \qquad B-A=\mathbb E\operatorname{Var}(Z\mid X).
 \label{eq:ci-analog-upper}
\end{equation}
At distinct times, conditionally independent readout noise gives
$\mathbb E[Z_0Z_t]=\langle f,e^{Qt}f\rangle_\pi$. Thus a valid upper confidence
bound on $B$ is also an upper bound on $A$. A lower bound on $B$ does not transfer.
Both $m$ and $B$ are unchanged under the additive reversible reference because
its stationary law and instantaneous detector channel are the same.

For example, let $Z=X/4+\epsilon/2$, with independent fair signs $X,\epsilon$.
Then $A=1/16$ while $B=5/16$. The observable second moment gives a useful upper
bound, but it is not the latent conditional second moment. Randomizing two
binary encodings from the same noisy $Z$ gives product mean $B$, not $A$, because
they share the analog measurement noise. Independently encoding readings at
distinct times is compatible with the correlation model.

The API accepts any valid simultaneous bounds on $m$ and $B$ and sets
$A\in[0,B_U]$. Its optional iid summary helper uses a two-sided Hoeffding bound
for the mean and a one-sided bound for the squared observation. A dependent
trajectory requires a separate valid calibration procedure; the iid helper
cannot be applied to every raw sample merely because the detector noise is
memoryless. Neither this upper-bound transfer nor the earlier exact transfer
is automatic for nonlinear finite-exposure functionals.

\section{A sharp calibrated reversible class}
For commensurate positive lags $k\Delta$, define
\begin{equation}
 \mathcal M_{m,A}=\left\{
 q_k=\frac{1+u+(A-u)\int_0^1x^k\,\mu(dx)}2:
 \mu\text{ a probability measure}\right\}.
 \label{eq:ci-calibrated-cone}
\end{equation}

\begin{theorem}[Sharp calibrated comparison class]
For a finite menu of positive integer lag indices, the closure of the agreement
vectors of finite stationary reversible chains with detector mean $m$ and
conditional second moment $A$ is exactly $\mathcal M_{m,A}$. The closure is
unchanged if the approximating chains must be irreducible.
\end{theorem}
\begin{proof}
Necessity follows from the established reversible spectral representation:
$C(t)=m^2+\sum_\ell b_\ell e^{-\lambda_\ell t}$, with
$b_\ell\geq0$ and $\sum_\ell b_\ell=A-m^2$~\cite{vlAldousFill2002}.
Additional conserved modes can be included at $\lambda=0$.

For sufficiency, a finite moment vector has an atomic representing measure
with at most $K+1$ atoms, where $K$ is the number of measured coordinates. If
$A=m^2$, use a constant detector. Otherwise put
\[
 a=\frac{m-A}{1-m},\qquad r=\frac{m-a}{1-a}.
\]
Then $-1\leq a<m<1$ and $0<r<1$. A two-state block with stationary law
$(r,1-r)$, detector values $(1,a)$, and generator
\[
 Q_x=\frac{-\log x}{\Delta}(\boldsymbol1\pi-I)
\]
has mean $m$, second moment $A$, and correlation
$m^2+(A-m^2)x^k$. Mix such blocks with the desired atomic weights. Replace any
atom $x=0$ by positive atoms approaching zero. Finally add
$\varepsilon(\boldsymbol1\pi-I)$ to the full generator. This preserves
$\pi,m,A$, makes the model irreducible and converges to the required moments as
$\varepsilon\downarrow0$.
\end{proof}

The sharpness statement concerns deterministic commensurate lags. A compatible
random-clock family preserves this calibrated set as an outer comparison class,
but its particular clock law can impose further restrictions.

\section{Calibration uncertainty and exact support}
A simultaneous interval for $m$ maps to an interval $[u_L,u_U]$ for $m^2$.
Its lower endpoint is zero if the mean interval crosses zero; otherwise use the
smaller squared endpoint. Together with $A\in[A_L,A_U]$, this produces the
exact polygon
\[
 \mathcal P=([u_L,u_U]\times[A_L,A_U])\cap\{0\leq u\leq A\leq1\}.
\]
For a contrast $c=(c_k)$, write $S_c=\sum_kc_k$ and
$p_c(x)=\sum_kc_kx^k$. Its reversible support is
\begin{equation}
 \sup_{q\in\mathcal M_{\mathcal P}}\sum_kc_kq_k
 =\frac12\max_{(u,A)\in\operatorname{vert}\mathcal P}
 \left[(1+u)S_c+(A-u)\max_{0\leq x\leq1}p_c(x)\right].
 \label{eq:ci-support}
\end{equation}
The full class is the convex hull of the vertex curves. A mixture has effective
parameters in $\mathcal P$, and reweighting its atomic measures by $A-u$ gives
another admissible measure. The implementation uses exact rational Bernstein
bounds for the continuous polynomial maximum; a numerical grid is only a lower
comparison.

For fixed $m$, these closed positive-lag classes are nested in $A$: extra
variance can be placed at $x=0$. Thus an upper calibration of $A$ sharpens the
comparison set. A lower calibration alone does not force visible variance when
arbitrarily fast hidden modes are allowed.

For $N$ independent same-state pair records, the count table gives
\[
 \widehat m=\frac{n_{++}-n_{--}}N,\qquad
 \widehat A=\frac{n_{++}+n_{--}-n_{+-}-n_{-+}}N.
\]
Both are averages of $N$ variables in $[-1,1]$. Hoeffding and a union bound give
a simultaneous radius
$\sqrt{2\log(4/\alpha)/N}$~\cite{caHoeffding1963}. This does not treat the
$2N$ readings as independent. Calibration error probability must be combined
with the subsequent inference guarantee. An empty physical calibration region
is reported explicitly rather than converted into an entropy estimate.

\section{Detecting excessively steep positive decay}
Since $x-x^2\leq1/4$, every calibrated reversible model satisfies the sharp bound
\begin{equation}
 q(t)-q(2t)\leq\frac{A-m^2}{8}.
 \label{eq:ci-steep-decay}
\end{equation}
A two-state block with $x=1/2$ attains it. This is a restriction on the rate of
change relative to the calibrated fluctuation amplitude, not merely on whether
the correlation is positive or decreasing.

Consider a synthetic four-state cycle with clockwise rate three, reverse rate
one and conditional detector means $(1/2,1/2,-1/2,-1/2)$. Its stationary law is
uniform, $m=0$, $A=1/4$, and
\[
 C(t)=\frac14 e^{-4t}\cos(2t).
\]
At $t=7/40$, exact series enclosures certify
\[
 q(t)-q(2t)>0.03473386533798
 >\frac1{32}=0.03125.
\]
Both correlations are positive and decrease. The uncalibrated support ceiling
$1/8$ does not detect this example. Its true generator entropy rate is
$2\log3$; that known model value is distinct from any lower bound inferred from
its restricted records. This is a population certificate, not sampled
experimental evidence.

\section{Exact positive-lag ambiguity}
\begin{theorem}[Finite pair menus do not identify the missing spectral mass]
For every integer $K\geq1$, there are finite irreducible reversible models with
zero detector mean, identical stationary binary-pair laws at each lag
$1,\ldots,K$, and different values of $A$.
\end{theorem}
\begin{proof}
Choose $x_j=j/(K+2)$ and $v_j=(-1)^j{K+1\choose j}$ for
$j=1,\ldots,K+1$. Finite differences give
\[
 \sum_{j=1}^{K+1}v_jx_j^k=0\quad(1\leq k\leq K),
 \qquad \sum_{j=1}^{K+1}v_j=-1.
\]
Split positive and negative coefficients and scale their magnitudes by
$s=(2^{K+1}-1)^{-2}$. The two positive spectral measures have equal measured
moments and unequal total mass.

For either measure, use independent symmetric telegraph coordinates $Z_j$ with
correlation $x_j^t$, and detector mean
$f(Z)=\sum_j\sqrt{s|v_j|}\,Z_j$. Its absolute value is at most one because the nonzero $|v_j|$ are integers,
so $\sqrt{s}\sum_j\sqrt{|v_j|}\leq\sqrt{s}\sum_j|v_j|\leq1$. The product chain is finite and irreducible, and its
correlation is exactly the desired spectral sum. Its mean is zero. The binary
pair law is $(1+abC(t))/4$ for outputs $a,b\in\{-1,1\}$, so the measured laws
agree exactly while $A=C(0)$ differs.
\end{proof}

For a small explicit example, compare spectral weights $3/16$ at $x=1/4$ and
$1/16$ at $x=3/4$ with weight $3/16$ at $x=1/2$. Their second moments are
$1/4$ and $3/16$, but both give $C(1)=3/32$ and $C(2)=3/64$. They have
irreducible four-state and two-state realizations, respectively, with exit rates
below one. Even a finite rate cap therefore does not imply exact identification
from an unchanged finite pair menu.

The result does not equate their full time-resolved trajectories or all
higher-order temporal records. Those observations, a known detector law or
true same-state calibration supply additional access. Repeating the unchanged
pair measurements cannot remove their exact ambiguity.

\section{Fast modes and the quantitative kinetic repair}
A second obstruction concerns uniform precision over unknown rates. Compare
independent fair binary outputs with a reversible two-state model having
correlation $Ae^{-\lambda t}$. Their pair total variation distance is
$Ae^{-\lambda t}/2$. For $N$ independent records at lags at least $t_{\min}>0$,
\begin{equation}
 \operatorname{TV}(P_0^{(N)},P_\lambda^{(N)})
 \leq\frac{NA}{2}e^{-\lambda t_{\min}}.
 \label{eq:ci-fast-mode}
\end{equation}
An order-one difference in $A$ can therefore be hidden at any finite precision
by a sufficiently fast reversible mode. Every model in this comparison has
zero stationary entropy production. For any finite independent clock menu with
no atom at zero, dominated convergence gives the same conclusion using
$\mathbb E e^{-\lambda T}$. Exact continuum data with unlimited precision are
a different observation contract.

A calibrated kinetic bound permits a controlled approximation. For any stationary
finite Markov chain, including a nonreversible one,
\begin{equation}
 A-C(t)=\frac12\mathbb E[(f(X_t)-f(X_0))^2].
\end{equation}
If every exit rate is at most $\Lambda$, then
\begin{equation}
 0\leq A-C(t)\leq2(1-e^{-\Lambda t}).
 \label{eq:ci-rate-repair}
\end{equation}
Indeed the integrand is at most four and vanishes when no jump occurs.
A bound $J$ on stationary jump activity similarly gives $A-C(t)\leq2Jt$.
These bounds make delayed calibration uniformly accurate as the accessible
delay approaches zero. They do not promise exact identification at fixed
positive lags. An energy-dissipation bound alone cannot supply either kinetic
bound: arbitrarily fast reversible chains dissipate no stationary heat.

Finally, the same-state noise premise is testable only with additional apparatus
knowledge. Conditionally anti-correlated readouts can each have means
$\pm1/2$ while their product has mean zero, although $A=1/4$. For example use
conditional count proportions $(1/2,1/4,1/4,0)$ and
$(0,1/4,1/4,1/2)$. An apparently feasible calibration can then be too small and
falsely reject an equilibrium record. Marginal counts alone do not establish
conditional independence or absence of measurement backaction.

The ingredients of this chapter---reversible spectral mixtures, truncated
moments, replicate covariance and concentration---are established. The scoped
advance is their integration into a sharp calibrated comparison class, an
executable confidence contract and explicit physical countermodels marking the
remaining access boundary. Historical originality and empirical validation
remain separate questions.
